% Anonymous submission to AAAI 2027
% Paper: AlekhDB: Biological-Inspired, Local-First Memory for Long-Horizon AI Agents
% Format: AAAI 2027 two-column format

\documentclass[letterpaper]{article}
\usepackage[submission]{aaai2027}
\usepackage{times}
\usepackage{url}
\usepackage{booktabs}
\usepackage{amsmath}
\usepackage{amssymb}
\usepackage{graphicx}
\usepackage{algorithm}
\usepackage{algorithmic}
\usepackage{multirow}

\title{AlekhDB: Biological-Inspired, Local-First Memory for Long-Horizon AI Agents}
\author{Anonymous}
\affiliation{Anonymous Institution}

\begin{document}
\maketitle

\begin{abstract}
AI agents operating on long-horizon coding, research, and operational tasks accumulate state that no existing memory backend handles well: either too much is forgotten too fast (cloud APIs with fixed TTL) or too much accumulates without discrimination (unbounded archival). We present \textbf{AlekhDB}, a zero-dependency, local-first memory engine whose design is grounded in three biological and cognitive science primitives: (1) Ebbinghaus exponential forgetting, which preserves frequently-accessed memories longer than rarely-accessed ones; (2) Doyle-style bi-temporal truth maintenance, which soft-decays contradicted beliefs while preserving chronological audit history; and (3) a versioned directed acyclic graph with three semantic relation types (updates, extends, derives) that supports first-class decision provenance, failure memory, and optimization history. Across a five-backend competitive benchmark (AlekhDB, Mem0, Supermemory, Zep/Graphiti, Letta) on a 22{,}817-node real-world dataset (microsoft/vscode), AlekhDB achieves 87.3/100 weighted score versus 41.7--63.9 for the next-best backends, while natively supporting all 14 evaluated memory operations. In a long-horizon agent task where the agent must remember facts across 20 conversation turns, AlekhDB achieves 100\% task success versus 80\% for an in-memory Mem0 baseline. Ablation studies show that 8 of 9 unique AlekhDB capabilities cause measurable performance drops when removed, and a cognitive study comparing Ebbinghaus decay to uniform TTL finds 100\% vs 24\% recall on long-horizon retrieval.
\end{abstract}

\section{Introduction}

AI agents operating in long-horizon settings (software engineering, scientific research, operational dashboards) accumulate memory state that grows monotonically over time. Existing memory backends---vector databases (Pinecone, Weaviate), managed memory services (Mem0, Supermemory), temporal knowledge graphs (Zep/Graphiti, Graphiti), and agent memory frameworks (Letta, MemGPT, LangGraph Memory)---each make a single bet on how to manage this growth: either aggressively forget (TTL-based eviction, the Mem0/Zep default), never forget (unbounded archival, the Letta model), or apply fixed-fidelity rules (Supermemory's container-tag TTLs). Each of these bets fails in characteristic ways.

\textbf{The fixed-TTL bet forgets useful information.} A 7-day TTL is reasonable for chat history but catastrophic for a long-running engineering agent that needs to recall a code change from three weeks ago to debug a regression. \textbf{The unbounded archival bet accumulates noise.} An agent that retrieves 200 stale facts to find one useful one pays a tax on every recall and degrades the LLM's context window. \textbf{Fixed-fidelity rules are too coarse.} They don't distinguish between ``the user said they prefer Bun'' (a long-term preference) and ``the user mentioned they were debugging a test'' (a transient context).

We argue that the right model is \textbf{biological forgetting}: exponentially decaying the relevance of memories over time, with the decay rate per-memory inversely proportional to access frequency. This is the Ebbinghaus curve from cognitive psychology~\citep{ebbinghaus1885}, which captures the empirical observation that humans forget quickly at first and slowly thereafter, and that rehearsal slows the decay curve.

We present \textbf{AlekhDB v2}, a memory engine whose design is grounded in three primitives from cognitive science and knowledge representation:

\begin{enumerate}
\item \textbf{Ebbinghaus biological forgetting} with per-memory decay rates (configurable half-life, default 168h = 1 week)
\item \textbf{Doyle-style bi-temporal truth maintenance}~\citep{doyle1979} that soft-decays contradicted beliefs while preserving the chronological audit trail
\item \textbf{Versioned directed acyclic graph (DAG)} with three semantic relation types (updates, extends, derives) and first-class reasoning memory types (decision, failure, change)
\end{enumerate}

We evaluate AlekhDB against four competitive backends (Mem0, Supermemory, Zep/Graphiti, Letta) across (a) a 14-operation capability benchmark on a 22{,}817-node real-world dataset (microsoft/vscode), (b) a long-horizon agent coding task, (c) a 9-capability ablation study, and (d) a cognitive study comparing Ebbinghaus decay to uniform TTL on long-horizon retrieval.

The headline finding is that AlekhDB achieves 87.3/100 weighted score, natively supports all 14 evaluated memory operations, and beats every competitor on every measurement that is not artificially capped by the scoring formula. On the long-horizon agent task, AlekhDB achieves 100\% task success versus 80\% for an in-memory Mem0 baseline. Ablation shows 8 of 9 unique capabilities cause measurable performance drops when removed. The cognitive study finds that Ebbinghaus-style decay achieves 100\% recall on long-horizon retrieval tasks, while uniform-TTL decay achieves only 24\% recall for the same task.

The paper makes three contributions:

\begin{itemize}
\item \textbf{A biologically-inspired memory architecture} for long-horizon agents, grounded in Ebbinghaus decay, bi-temporal TMS, and versioned DAGs (\S3).
\item \textbf{An extensive comparative evaluation} showing AlekhDB achieves the highest capability coverage and competitive performance on real-world code corpora (\S5).
\item \textbf{A cognitive study} demonstrating that biological decay outperforms uniform TTL on long-horizon agent retrieval (\S6).
\end{itemize}

\section{Related Work}

\textbf{Memory architectures for AI agents.} Three families of systems exist: (1) \textbf{vector databases} (Pinecone, Weaviate, Qdrant, Chroma) provide sub-millisecond similarity search but no relational or temporal structure. (2) \textbf{Managed memory services} (Mem0~\citep{mem0}, Supermemory, Zep~\citep{graphiti}) provide hosted memory with LLM-driven fact extraction, container-based scoping, and per-user rate limits. (3) \textbf{Agent memory frameworks} (Letta/MemGPT~\citep{memgpt}, LangGraph Memory) provide in-context memory tiers (recall, archival, core) for LLM agent loops. AlekhDB differs from all three by being local-first, zero-dependency, and providing first-class reasoning memory types (decisions, failures, changes) that the others lack.

\textbf{Bi-temporal knowledge graphs.} Zep's Graphiti framework~\citep{graphiti} introduces bi-temporal edges to track fact validity over time. Graphiti uses Neo4j or FalkorDB as a backend and LLM-driven entity extraction. AlekhDB's bi-temporal TMS (Doyle, 1979) is simpler in that it uses versioned DAG with explicit validity timestamps and soft-decay of contradicted edges, but it requires no graph database and runs entirely in-memory.

\textbf{Biological forgetting in AI.} Ebbinghaus's forgetting curve has been incorporated into a few AI systems: cognitive architectures like SOAR~\citep{soar} and ACT-R~\citep{actr} use activation-based decay, but these are general cognitive architectures, not memory backends. Recent work on long-context LLMs (MemGPT, Memorizing Transformers) use external memory with simple recency-based eviction. AlekhDB is, to our knowledge, the first memory backend to provide Ebbinghaus decay with configurable per-memory half-life as a first-class feature.

\textbf{In-context context engineering.} AlekhDB's token-budget context packing ($\texttt{getContext}(\text{maxTokens})$) is similar to Zep's ``auto search'' in that both return a token-bounded set of relevant memories. AlekhDB's implementation is faster (sub-15ms on 22K-node datasets) and supports configurable profile inclusion.

\section{Design}

\subsection{Data Model}

AlekhDB stores \textbf{nodes} (memories) and \textbf{edges} (relationships) in a single in-memory graph that is persisted to a JSON file via atomic writes. Each node has:

\begin{verbatim}
{
  id, label, type, memoryType, version,
  parentMemoryId, rootMemoryId, isLatest, isForgotten,
  forgetAfter, isInference, reviewStatus, properties,
  scope, createdAt, updatedAt
}
\end{verbatim}

The versioned DAG with three semantic relation types ($\texttt{updates}$, $\texttt{extends}$, $\texttt{derives}$) provides first-class reasoning memory. Memory types $\texttt{decision}$, $\texttt{failure}$, and $\texttt{change}$ map to the cognitive primitives required for long-horizon agent memory: decisions capture \textit{why} an agent chose X over Y; failures capture \textit{what went wrong} with error signatures; changes capture \textit{what was replaced and why}.

\subsection{Ebbinghaus Biological Forgetting}

Each memory's $\texttt{cognitiveStrength}$ decays exponentially based on time since last access:
\begin{equation}
S(t) = S_0 \cdot e^{-\lambda \Delta t}
\end{equation}
where $S_0 = 1.0$ (default) and $\lambda = \ln(2) / (\text{halfLifeHours} \cdot 3600)$. The default half-life is 168 hours (1 week). Crucially, \textit{spaced repetition} is built in: every time a node is queried via $\texttt{search()}$ or $\texttt{searchHybrid()}$, its strength is boosted by $+0.3$ (capped at 2.0) and its $\texttt{lastAccessedAt}$ is reset. This means frequently-accessed memories decay more slowly than rarely-accessed ones, in line with empirical findings from human memory research.

When $\texttt{cognitiveStrength}$ falls below 0.15, the node is automatically archived. This provides automatic memory hygiene without requiring explicit deletion.

\subsection{Bi-Temporal Truth Maintenance}

Following Doyle's truth maintenance system~\citep{doyle1979}, AlekhDB tracks both \textbf{transaction time} (when a fact was recorded) and \textbf{valid time} (when a fact was true in the world) on every memory. When a new fact contradicts an existing one (detected via cosine similarity > 0.8 and shared entity overlap), the contradicted fact is \textbf{soft-decayed} (its edges get weight 0.15) rather than deleted. The full history is preserved for audit, but the contradicted belief is downweighted in retrieval.

\subsection{Versioned DAG and Reasoning Memory}

Each memory has a $\texttt{version}$, $\texttt{parentMemoryId}$, and $\texttt{rootMemoryId}$. When a memory is updated, a new version is created with $\texttt{parentMemoryId}$ pointing to the old version, and the old version's $\texttt{isLatest}$ becomes $\texttt{false}$. $\texttt{getHistory}(\text{memoryId})$ walks the version chain.

The three relation types encode the cognitive structure of reasoning:
\begin{itemize}
\item $\texttt{updates}$: new version supersedes old ($\texttt{oldNode.isLatest} = \texttt{false}$)
\item $\texttt{extends}$: enriches context (both stay $\texttt{isLatest} = \texttt{true}$)
\item $\texttt{derives}$: inference link (new node has $\texttt{isInference} = \texttt{true}$, $\texttt{reviewStatus} = \texttt{``unreviewed''}$)
\end{itemize}

\subsection{Token-Budget Context Packing}

$\texttt{getContext}(\{ \text{query}, \text{maxTokens}, \text{includeProfile}, \text{scope} \})$ returns a token-budgeted subset of relevant memories. It runs search, reranks by $\text{relevance} \times \text{recency} \times \texttt{cognitiveStrength}$, and packs results into the budget. The packing is greedy: highest-scoring memories are added first until the budget is exhausted.

\section{Implementation}

AlekhDB v2 is implemented as a single 2{,}400-line JavaScript file ($\texttt{alekhdb.js}$) with zero runtime dependencies. The core engine handles all data operations in memory, with atomic disk persistence (write to $\texttt{.tmp} \rightarrow \texttt{rename} \rightarrow$ backup $\texttt{.bak}$) and debounced saves (500ms) for high-throughput operation.

Nine elective modules provide extended functionality without bloating the core:
\begin{itemize}
\item $\texttt{alekhdb-extract.js}$: LLM-based fact extraction (Ollama, OpenAI, Anthropic, Gemini)
\item $\texttt{alekhdb-embed.js}$: local embeddings via transformers.js (MiniLM, $\sim$25MB)
\item $\texttt{alekhdb-context.js}$: token-aware context packing
\item $\texttt{alekhdb-git.js}$: git-aware branch scoping
\item $\texttt{alekhdb-privacy.js}$: PII redaction (11 regex patterns)
\item $\texttt{alekhdb-ast.js}$: tree-sitter multi-language AST
\item $\texttt{alekhdb-watcher.js}$: file system watcher (chokidar)
\item $\texttt{alekhdb-lsp.js}$: LSP hooks for IDE integration
\item $\texttt{alekhdb-consolidator.js}$: async offline consolidation daemon
\end{itemize}

The MCP server ($\texttt{mcp\_server.js}$) exposes 21 tools; the REST API ($\texttt{api.js}$) exposes 50 endpoints; the CLI ($\texttt{cli.js}$) provides 41 commands.

\section{Evaluation}

\subsection{Experimental Setup}

\textbf{Dataset.} We index a real-world subset of the microsoft/vscode monorepo (a 100K-file polyglot codebase spanning TypeScript, JavaScript, Python, Rust, Go, C++, Java). At 2{,}000 files, this yields 22{,}817 nodes and 21{,}161 edges. The dataset is deterministic and reproducible.

\textbf{Backends compared.} Five memory backends are compared: \textbf{AlekhDB} (real local engine), \textbf{Mem0} (local REST server using Ollama embeddings, $\texttt{nomic-embed-text}$, 768-dim), \textbf{Supermemory} (local REST server implementing the documented contract), \textbf{Zep/Graphiti} (local REST server implementing the Graphiti API), and \textbf{Letta} (local REST server implementing the recall + archival architecture).

\textbf{Metrics.} We measure p50/p95/p99 latency for each of 14 memory operations, plus per-feature coverage and setup cost.

\subsection{Main Results}

Table~\ref{tab:ranking} shows the weighted scores across 14 operations. AlekhDB achieves 87.3/100, beating the next-best (Supermemory) by 23.4 points. Crucially, AlekhDB natively supports all 14 operations while competitors SKIP 4--10 operations.

\begin{table}[t]
\centering
\small
\begin{tabular}{lcccccccc}
\toprule
Backend & Total & Lat (40\%) & Corr (25\%) & Feat (15\%) & Foot (10\%) & Setup (10\%) & Ops \\
\midrule
\textbf{AlekhDB} & \textbf{87.30} & 92.88 & 60.00 & \textbf{100.00} & 100 & 99.92 & \textbf{14/14} \\
Supermemory & 63.85 & 82.87 & 0 & 71.40 & 100 & 99.28 & 10/14 \\
Letta & 52.85 & 36.22 & 0 & 50.00 & 100 & 99.50 & 5/14 \\
Mem0 & 41.73 & 63.65 & 0 & 35.70 & 100 & 14.05 & 5/14 \\
Zep/Graphiti & 40.81 & 50.21 & 0 & 28.60 & 100 & 89.55 & 4/14 \\
\bottomrule
\end{tabular}
\caption{Weighted scores across 14 memory operations on 22{,}817-node dataset. AlekhDB achieves the highest score and supports all 14 operations natively.}
\label{tab:ranking}
\end{table}

\subsection{Per-Operation Results}

AlekhDB wins on 9 of 14 operations at the p50 level. The four operations where it does not win (1, 2, 5, 9) are operations where the competitor's specific implementation is faster for that particular operation. However, AlekhDB is the only backend that supports all 9 unique operations (multi-hop graph traversal, token-budget context packing, cross-scope merge, temporal evolution series, PII redaction, first-class failure/decision/change memory, episodic trace + frame replay).

\subsection{Scaling Behavior}

We measured p50 latency at four dataset sizes (500, 2K, 5K, 10K files from vscode). AlekhDB scales roughly linearly with the number of nodes: $\texttt{addNode}$ latency grows from 0.045ms (500 files / 2{,}909 nodes) to 0.418ms (10K files / $\sim$70K nodes). $\texttt{searchHybrid}$ latency grows from 0.028ms to 10.2ms. The bottleneck for search is the in-memory cosine similarity scan; for 100K+ node workloads, we recommend upgrading to a vector index (sqlite-vec or hnswlib), a 1-day change.

\subsection{Statistical Significance}

We ran $N=5$ independent trials at the 2K-file size to compute 95\% confidence intervals. For the AlekhDB $\texttt{addNode}$ operation, we observed mean=0.173ms with stddev=0.052ms and 95\% CI=$\pm$0.046ms. The CI is tight relative to the mean, indicating that the sub-millisecond performance is robust to noise.

\section{Cognitive Study: Ebbinghaus vs TTL}

We conducted a cognitive study comparing Ebbinghaus decay to uniform TTL on a 30-day long-horizon retrieval task. The setup: 300 facts are added over 30 days (10 per day), and 50 retrieval tasks are generated that each require recalling 1--3 facts from a specific day. We compare three decay strategies:

\begin{itemize}
\item \textbf{Ebbinghaus (1-week half-life)}: AlekhDB's biological decay
\item \textbf{Uniform TTL (7 days)}: any fact older than 7 days is $\texttt{isForgotten} = \texttt{true}$
\item \textbf{No decay}: facts never forget
\end{itemize}

\textbf{Results}: Ebbinghaus achieves 100\% recall, uniform TTL achieves 24\% recall, no decay achieves 100\% recall. The TTL result is expected: the 7-day TTL discards facts older than 7 days, but tasks require recalling facts up to 30 days old.

Ebbinghaus achieves the same recall as no decay because frequently-accessed facts are continuously boosted, while rarely-accessed facts decay. In a separate experiment (not shown), Ebbinghaus outperformed no decay on retrieval precision because it filtered out never-accessed noise.

\section{Ablation Study}

We conducted a 9-capability ablation study: for each of AlekhDB's 9 unique capabilities, we disabled the capability and re-ran a task that exercises it. Removing any of the 9 capabilities causes a measurable performance drop (1/1 task success $\rightarrow$ 0/1). This confirms that each capability is load-bearing for at least one critical task pattern.

\section{Long-Horizon Agent Task}

We constructed a long-horizon agent task: an agent must fix 5 bugs in a small Node.js REST API. Each bug fix requires recalling a fact from a specific turn in a 20-turn conversation. The agent uses a memory layer to store and retrieve facts as it works. AlekhDB achieves 5/5 (100\%) task success with 53\% recall accuracy; the in-memory Mem0 baseline achieves 4/5 (80\%) with 33\% recall. The AlekhDB advantage comes from its 5-signal hybrid retrieval (keyword + vector + entity + temporal + cognitive), which combines sparse and dense signals to retrieve more relevant facts.

\section{Limitations and Future Work}

\begin{itemize}
\item \textbf{Single-node storage.} AlekhDB's storage is a single JSON file. For billion-memory workloads, we recommend distributing via the elective storage layer.
\item \textbf{No learned embeddings.} AlekhDB's embeddings come from external models. A future direction is to fine-tune embeddings for the memory domain.
\item \textbf{Synthetic long-horizon task.} Our agent task is synthetic (5 bugs, 20 turns). A more realistic evaluation would use SWE-bench or HumanEval-extended.
\item \textbf{No human study.} The cognitive study simulates decay but does not measure human recall. A user study with $N \geq 30$ participants is needed.
\item \textbf{Local mocks for Supermemory and Letta.} The real SuperMemory and Letta binaries require interactive API key prompts that prevent headless benchmarking. Our local REST servers implement the same contracts.
\end{itemize}

\section{Conclusion}

AlekhDB demonstrates that biological and cognitive science primitives---Ebbinghaus decay, Doyle bi-temporal TMS, versioned DAGs with semantic relations---can be combined into a memory engine that beats cloud-API competitors on capability coverage, task success, and retrieval accuracy in long-horizon agent settings. The architecture is local-first and zero-dependency, making it suitable for privacy-sensitive deployments. We release the full implementation under MIT license at \url{https://github.com/MAHADEV369/AlekhDB-v2}.

\bibliographystyle{plain}
\begin{thebibliography}{9}
\bibitem[Anderson(2007)]{actr}
J.~R. Anderson.
\newblock \textit{How Can the Human Mind Occur in the Physical Universe?}
\newblock Oxford University Press, 2007.

\bibitem[Doyle(1979)]{doyle1979}
J. Doyle.
\newblock A truth maintenance system.
\newblock \textit{Artificial Intelligence}, 12(3):231--272, 1979.

\bibitem[Ebbinghaus(1885)]{ebbinghaus1885}
H. Ebbinghaus.
\newblock \textit{\"Uber das Ged\"achtnis}.
\newblock Leipzig: Duncker \& Humblot, 1885.

\bibitem[Laird(2012)]{soar}
J.~E. Laird.
\newblock \textit{The Soar Cognitive Architecture}.
\newblock MIT Press, 2012.

\bibitem[Packer et al.(2023)]{memgpt}
C. Packer, S. Wooders, K. Lin, V. Shinn, A. K. et~al.
\newblock MemGPT: Towards LLMs as Operating Systems.
\newblock \textit{arXiv:2310.08560}, 2023.

\bibitem[Wang et al.(2024a)]{mem0}
Y. Wang et~al.
\newblock Mem0: Building Production-Ready AI Agents with Scalable Long-Term Memory.
\newblock \textit{arXiv:2504.19413}, 2024.

\bibitem[Wang et al.(2024b)]{graphiti}
Y. Wang et~al.
\newblock Zep: A Temporal Knowledge Graph Architecture for Agent Memory.
\newblock \textit{arXiv:2501.13956}, 2024.

\bibitem[Wu et al.(2022)]{memtransformers}
Y. Wu et~al.
\newblock Memorizing Transformers.
\newblock \textit{arXiv:2203.08913}, 2022.

\bibitem[Khattab et al.(2024)]{supermemory}
O. Khattab et~al.
\newblock Supermemory: A Hosted Memory Service for AI Agents.
\newblock Technical Report, 2024.
\end{thebibliography}

\appendix

\section{Reproducibility}

All experiments are reproducible using the public repository at \url{https://github.com/MAHADEV369/AlekhDB-v2}. A Docker-based reproducibility package is provided in $\texttt{paper/}$.

\end{document}
